\section{USAMO 1998}
        \begin{enumerate}
        	\item (\textit{USAMO 1998}) Suppose that the set $\{1,2,\cdots, 1998\}$ has been partitioned into disjoint pairs $\{a_i,b_i\}$ ($1\leq i\leq 999$) so that for all $i$, $|a_i-b_i|$ equals $1$ or $6$. Prove that the sum

 \[|a_1-b_1|+|a_2-b_2|+\cdots +|a_{999}-b_{999}|\]
ends in the digit $9$.


 

	\item (\textit{USAMO 1998}) Let ${\cal C}_1$ and ${\cal C}_2$ be  concentric circles, with ${\cal C}_2$ in the interior of  ${\cal C}_1$. From a point $A$ on ${\cal C}_1$ one draws the tangent $AB$ to ${\cal C}_2$ ($B\in {\cal C}_2$). Let $C$ be the second point of intersection of $AB$ and ${\cal C}_1$, and let $D$ be the midpoint of $AB$. A line passing through $A$ intersects ${\cal C}_2$ at $E$ and $F$ in such a way that the perpendicular  bisectors of $DE$ and $CF$ intersect at a point $M$ on $AB$. Find, with proof,  the ratio $AM/MC$.


 

	\item (\textit{USAMO 1998}) Let $a_0,a_1,\cdots ,a_n$ be numbers from the interval $(0,\pi/2)$ such that

 \[\tan \left(a_0-\frac{\pi}{4}\right)+ \tan \left(a_1-\frac{\pi}{4}\right)+\cdots +\tan \left(a_n-\frac{\pi}{4}\right)\geq n-1.\]
Prove that

 \[\tan a_0\tan a_1 \cdots \tan a_n\geq n^{n+1}.\]


 

	\item (\textit{USAMO 1998}) A computer screen shows a $98 \times 98$ chessboard, colored in the usual way. One can select with a mouse any rectangle with sides on the lines of the chessboard and click the mouse button: as a result,  the colors in the selected rectangle switch (black becomes white, white becomes black). Find, with proof, the minimum number of mouse clicks needed to make the chessboard all one color.


 

	\item (\textit{USAMO 1998}) Prove that for each $n\geq 2$, there is a set $S$ of $n$ integers such that $(a-b)^2$ divides $ab$ for every distinct $a,b\in S$.


 

	\item (\textit{USAMO 1998}) Let $n \geq 5$ be an integer. Find the largest integer $k$ (as a function of $n$) such that there exists a convex $n$-gon $A_{1}A_{2}\dots A_{n}$ for which exactly $k$ of the quadrilaterals $A_{i}A_{i+1}A_{i+2}A_{i+3}$ have an inscribed circle. (Here $A_{n+j} = A_{j}$.)


 

\end{enumerate}